\documentclass{article}
\usepackage{hyperref}
\usepackage{geometry}
\geometry{a4paper, margin=1.5cm}
\pagestyle{empty}

\title{Étude des Polynômes du Second Degré}
\author{}
\date{}

\begin{document}

\begin{center}
\section*{Activité 2: Étude des polynômes du second degré}
\end{center}
\bigskip
\bigskip

\noindent
\textbf{Instructions :} \\
Ouvrez GeoGebra et entrez l'équation suivante dans la barre d'entrée :
\[
f(x) = ax^2 + bx + c
\]
ou bien cliquez sur le lien suivant pour ouvrir une version préparée de l'activité :
\href{https://www.geogebra.org/m/swbegdnu}{https://www.geogebra.org/m/swbegdnu}

\bigskip
Une fois l'application ouverte, vous pourrez modifier les valeurs des paramètres $a$, $b$, et $c$ et observer leurs effets sur la courbe du polynôme.

\bigskip
\noindent
\textbf{Énoncé :} \\
Nous allons étudier la forme générale d'un polynôme du second degré $f(x) = ax^2 + bx + c$ et explorer comment les paramètres $a$, $b$ et $c$ influencent la courbe, les racines et le discriminant $\Delta$.

\begin{enumerate}
    \item On fixe $a = 1$, $b = 0$, et $c = -1$.
    \begin{enumerate}
        \item Faites varier $c$. Qu'observez-vous ?
        \item Pour quelle valeur de $y$ la courbe coupe-t-elle l'axe des ordonnées ?
        \item Cherchons les points où la courbe coupe l'axe des abscisses.
        \begin{enumerate}
            \item Mettez $c = -4$ et trouvez graphiquement les valeurs de $x$ pour lesquelles la courbe coupe l'axe des abscisses.
            \item À quelle équation cela correspond-il ?
            \item Quand $c < 0$, quelles sont les solutions de cette équation ?
            \item Pour quelle valeur de $c$ a-t-on $x_1 = x_2$ ?
            \item Pour quelles valeurs de $c$ l'équation n'admet-elle pas de solutions ?
        \end{enumerate}
    \end{enumerate}

    \item On fixe $a = 1$, $b = 0$, et $c = -4$.
    \begin{enumerate}
        \item Faites varier $a$. Qu'observez-vous ?
        \item Pour quelle valeur de $y$ la courbe coupe-t-elle l'axe des ordonnées ?
        \item Cherchons les points où la courbe coupe l'axe des abscisses.
        \begin{enumerate}
            \item À quelle équation cela correspond-il ?
            \item Pour quelles valeurs de $c$ et de $a$ cette équation a-t-elle des solutions ?
        \end{enumerate}
    \end{enumerate}

    \item On fixe $a = -0.5$.
    \begin{enumerate}
        \item Faites varier $b$. Qu'observez-vous ?
        \item À quoi correspond la valeur de $b$ dans ce cas ?
        \item Testez pour d'autres valeurs de $a$. Qu'observez-vous ?
    \end{enumerate}

    \item Soit $a,b,c$ quelconque avec $a \neq 0$.
    Cherchons les points où la courbe coupe l'axe des abscisses.
        \begin{enumerate}
            \item À quelle équation cela correspond-il ?
            \item En complétant le carré, reformulez cette équation.
            \item Pour quelles valeurs de $a$, $b$, et $c$ l'équation a-t-elle des solutions ?
            \item Posez $\Delta = b^2 - 4ac$. Quelles sont alors les solutions de l'équation ?
        \end{enumerate}
\end{enumerate}

\newpage


\end{document}